% !TeX root = ../ejemplo2.tex

\section*{Organización del documento}

El presente documento se encuentra estructurado por los siguientes capítulos y anexos:

En el capítulo 1, se encuentran los antecedentes que permitieron identificar la problemática que rige este proyecto, lo que a su vez nos guío al análisis de esta misma para sustentar la propuesta de solución planteada a lo largo de este documento \\

En el capítulo 2, se presenta la definición del proyecto, donde se explica la justificación desde un contexto social, personal y académico del proyecto y además, se da una explicación breve y concisa sobre el alcance, los requerimientos, la arquitectura, las propiedades de software, el modelo de información y la metodología de desarrollo. \\

En el capítulo 3, se encuentra el estado del arte, en donde se revisaron y analizaron trabajos anteriores con similitudes al proyecto a realizar, esto con el fin de realizar una comparación de las principales características e identificar posibles áreas de mejora.\\

En el capítulo 4, se presenta el marco teórico donde se explican los términos, descripciones y demás elementos importantes que se relacionan en menor o mayor medida con este trabajo a fin de tener un panorama general sobre la terminología utilizada.\\

En el capítulo 5, se presenta los detalles del trabajo realizado, es decir se explica cómo y con que tecnologías y métodos se implementó cada una de las secciones del sistema.\\

En el capítulo 6, se exponen las pruebas realizadas en la aplicación móvil a los casos de uso más importantes, para buscar errores para su posterior corrección y/o determinar si la funcionalidad y/o proceso se está realizando de manera satisfactoria o cumple con lo esperado, además, se muestran las pruebas hechas al sistema de reconocimiento facial.\\

En el capítulo 7, se presentan los resultados obtenidos de las encuestas hechas a los alumnos, docentes y personal de seguridad de la escuela superior de cómputo (ESCOM), sobre el funcionamiento de la aplicación, las opiniones sobre esta y si les parece útil o si creen que podría ayudar reducir el problema de la suplantación de identidad durante los ETS. También, se presentan los resultados de las pruebas de la red neuronal y sus matrices de confusión.\\

Finalmente, en el capítulo 8, se presentan las direcciones futuras y consideraciones identificadas para este proyecto, una vez culminada la fase de Trabajo Terminal actual. Si bien se ha logrado la implementación de la aplicación móvil y el sistema de gestión asociado, incluyendo el reconocimiento facial en tiempo real, este capítulo delinea las potenciales expansiones y mejoras que podrían llevarse a cabo en una eventual continuación del proyecto, más allá de los objetivos académicos de este Trabajo Terminal.\\

En el anexo A, se sustenta la metodología de trabajo empleada, se presentan los resultados del estudio de factibilidad realizado y el análisis del sistema que compone la propuesta de solución a fin de delimitar el alcance de este y los elementos a desarrollar, asimismo, se presenta el análisis de riesgos correspondiente.\\  

En el Anexo B, se muestran las decisiones de diseño tomadas para este proyecto, esto se encuentra representado a través de diagramas para ofrecer una visualización y explicación del funcionamiento del sistema y las interacciones que existen entre los componentes que lo conforman. Estos incluyen el diagrama de estructura general del sistema, la estructura interna de los módulos, diagramas de clases, diagramas de componentes y diagramas de secuencia.\\



% Finalmente en el capítulo 6 se encuentra en trabajo futuro identificado para este trabajo terminal 1, destacando la implementación de la aplicación móvil, el sistema para la gestión de alumnos, personal y exámenes ETS así como la implementación del sistema de reconocimiento facial en tiempo real. Se espera que este trabajo se desarrolle y culmine durante la segunda etapa del trabajo terminal.